\section{Analyse et choix de l'algorithme}

\subsection{Résultats et complexité}
Les tests confirment la complexité théorique en O(NlogN).
Quand on multiplie la taille des données par 10,
le temps de calcul augmente de façon cohérente (environ 15 à 20 fois).

On note cependant que sur 2,5 millions de lignes,
le MergeSort est 3,5 fois plus rapide que le QuickSort (0,40s contre 1,41s).
Cela s'explique par la sensibilité du QuickSort au choix du pivot dans une liste déjà triée,
alors que le MergeSort reste stable et constant peu importe l'ordre des données.

\subsection{Algorithme approprié}
Le tri fusion (MergeSort) est préféré.

Même si le QuickSort est un peu plus rapide,
le MergeSort garantit une performance stable en O(NlogN) dans tous les cas.
Nos tests prouvent que le MergeSort est plus performant sur de gros volumes de données déjà triées.

Il faut noter que les performances du tri rapide dépendent du choix du pivot.
Avec un autre pivot(nous avons utilisé l'élément le plus à droite), il peut être plus efficace.
